\subsection{Dados de entrada}

O dataset contém 8 variáveis preditoras e 1 variável alvo, que é o salário bruto.

\begin{table}[htbp]
    \centering
    \caption{Variáveis preditoras}
    \begin{tabular}{|l|l|l|}
        \hline
        \textbf{Variável} & \textbf{Descrição} & \textbf{Tipo} \\ \hline
            year & Ano do registro & Numérica \\ \hline
            age & Idade do trabalhador & Numérica \\ \hline
            maritl & Estado civil & Categórica \\ \hline
            race & Raça & Categórica \\ \hline
            education & Educação & Categórica \\ \hline
            jobclass & Classe de emprego (Industrial ou Informática) & Categórica \\ \hline
            health & Nível de saúde (bom ou muito bom) & Categórica \\ \hline
            health\_ins & Possui plano de saúde (sim ou não) & Categórica \\ \hline
    \end{tabular}
\end{table}

O conjunto possui 3.000 observações e não apresenta valores ausentes. A coluna \texttt{Unnamed: 0}, presente no arquivo original, representa apenas o índice das observações e será descartada antes da modelagem. Após essa remoção, permanecem as oito variáveis preditoras e a variável alvo \texttt{wage}.

\begin{table}[htbp]
    \centering
    \caption{Resumo da variável alvo \texttt{wage}}
    \begin{tabular}{lr}
        \toprule
        \textbf{Estatística} & \textbf{Valor} \\
        \midrule
        Número de observações & 3.000 \\
        Valores ausentes & 0 \\
        Média & 111,70 \\
        Mediana & 104,92 \\
        Mínimo & 20,09 \\
        Máximo & 318,34 \\
        Observações além de $1{,}5\,\mathrm{IQR}$ & 110 \\
        Observações além de $3\,\mathrm{IQR}$ & 77 \\
        \bottomrule
    \end{tabular}
\end{table}

A média superior à mediana indica assimetria à direita na distribuição salarial. Entre as 110 observações situadas além de $1{,}5\,\mathrm{IQR}$, 109 estão acima do limite superior e uma está abaixo do limite inferior. As 77 observações além de $3\,\mathrm{IQR}$ estão todas no extremo superior da distribuição. Essas observações serão mantidas, pois podem representar salários legítimos e sua remoção alteraria o problema de previsão definido pelo estudo.

As variáveis \texttt{year} e \texttt{age} serão tratadas como numéricas e normalizadas para o intervalo entre 0 e 1 pelo método Min-Max. As seis variáveis restantes serão tratadas como categóricas e convertidas em variáveis indicadoras por codificação \textit{one-hot}, com a primeira categoria de cada variável utilizada como referência. Após essas transformações, o conjunto de preditores será composto por 16 colunas.

Após o pré-processamento, os dados serão divididos aleatoriamente em conjuntos de treinamento e teste, contendo, respectivamente, 80\% e 20\% das observações. Dessa forma, serão destinadas aproximadamente 2.400 observações ao treinamento e 600 ao teste. Será utilizada a semente aleatória 42 para permitir a reprodução da divisão, e o conjunto de teste será reservado para a avaliação final dos modelos.

\subsection{Técnicas utilizadas}

Optamos por comparar três abordagens de regressão linear para a previsão do salário bruto: regressão linear simples, regressão linear múltipla e regressão Lasso. Essa delimitação permite avaliar o ganho obtido pela combinação de variáveis e verificar se a regularização melhora o desempenho preditivo.

\textbf{Regressão Linear.} A regressão linear simples estima o salário a partir de uma única variável preditora. Para cada uma das oito variáveis disponíveis, será ajustado um modelo independente, permitindo avaliar isoladamente sua capacidade preditiva. A previsão é obtida pela combinação linear

\begin{equation}
    \widehat{y}_i = \beta_0 + \beta_1 x_i,
\end{equation}

\noindent em que $\widehat{y}_i$ é o salário previsto para o indivíduo $i$, $x_i$ é o valor da variável preditora, $\beta_0$ é o intercepto e $\beta_1$ é o coeficiente associado à variável. Os coeficientes são estimados de modo a minimizar a soma dos quadrados dos resíduos entre os valores observados e previstos. Quando a variável preditora for categórica, suas categorias serão representadas por variáveis indicadoras.

\textbf{Regressão Linear Múltipla.} A regressão linear múltipla estende o modelo anterior ao empregar simultaneamente duas ou mais variáveis preditoras. Sua combinação linear é dada por

\begin{equation}
    \widehat{y}_i = \beta_0 + \sum_{j=1}^{p} \beta_j x_{ij}
    = \beta_0 + \beta_1 x_{i1} + \cdots + \beta_p x_{ip},
\end{equation}

\noindent em que $p$ representa o número de variáveis utilizadas, $x_{ij}$ é o valor da variável $j$ para o indivíduo $i$ e $\beta_j$ é seu respectivo coeficiente. Esse modelo permite considerar conjuntamente as informações disponíveis, embora pressuponha uma relação linear entre os preditores e o salário. As variáveis categóricas serão previamente representadas por variáveis indicadoras.

\textbf{Regressão Lasso.} O Lasso estende a regressão linear múltipla ao acrescentar uma penalização $L_1$ à função de ajuste. Os coeficientes são estimados pela minimização de

\begin{equation}
    \sum_{i=1}^{n}\left(y_i-\widehat{y}_i\right)^2
    + \alpha\sum_{j=1}^{p}\left|\beta_j\right|,
\end{equation}

\noindent em que $\alpha$ controla a intensidade da regularização. Valores pequenos de $\alpha$ produzem resultados próximos aos da regressão linear múltipla, enquanto valores maiores reduzem os coeficientes e podem fazer com que alguns deles se tornem exatamente zero. Dessa forma, o Lasso também realiza seleção de características. Para que a penalização seja aplicada de forma comparável, serão utilizadas as variáveis previamente codificadas e normalizadas.

\subsection{Métricas}

O desempenho preditivo dos modelos será avaliado pelo Erro Absoluto Médio (MAE), definido por

\begin{equation}
    \mathrm{MAE} = \frac{1}{n}\sum_{i=1}^{n}\left|y_i - \widehat{y}_i\right|.
\end{equation}

\noindent Essa métrica representa a magnitude média dos erros na mesma unidade da variável salarial e possui interpretação direta. Além disso, por não elevar os resíduos ao quadrado, é menos influenciada por salários extremos. Essa característica é relevante para os dados utilizados, cuja distribuição salarial apresenta assimetria à direita. Assim, o MAE será adotado como critério para a seleção do melhor modelo: quanto menor seu valor, melhor o desempenho preditivo.

O MAE será calculado sobre observações não utilizadas no ajuste dos modelos. Na comparação por validação cruzada, serão empregadas as mesmas divisões dos dados para todas as técnicas, reportando-se a média e o desvio-padrão do MAE. Como referência, o desempenho também será comparado ao de um modelo ingênuo que prevê sempre a mediana salarial do conjunto de treinamento.

\subsection{Parâmetros testados}

As regressões lineares simples e múltipla não possuem hiperparâmetros. Serão ajustados oito modelos simples, um para cada variável preditora, e um modelo múltiplo contendo as oito variáveis.

No Lasso, serão testados 200 valores do parâmetro de regularização $\alpha$, distribuídos uniformemente em escala logarítmica entre $10^{-4}$ e $10^{2}$. O intervalo abrange desde penalizações quase nulas até penalizações capazes de reduzir diversos coeficientes a zero.

Cada configuração será avaliada por validação cruzada com cinco divisões, sempre com as mesmas partições, e a seleção será baseada no menor MAE médio. Ao todo, serão avaliadas 209 configurações distintas e realizados 1.045 ajustes durante a validação cruzada.

Após a seleção, a melhor configuração será reajustada com todo o conjunto de treinamento e avaliada uma única vez no conjunto de teste.

\newpage
